\scanpage{257}
\mspage{131}
% archive sheet number: 256

La condition sur les sommes assure que dans $\mathrm{Hot}_W$, la somme de deux objets existe, et que $\mathrm{Cat}\to\mathrm{Hot}_W$ commute aux sommes de deux objets. Ceci dit, soit
\[ f:X\longrightarrow Y \]
une flèche telle que $f\in W$, mais $\pi_0(f)$ non surjectif. \inserted{Supposons d'abord $X\ne\varnothing$ (cas le moins évident).}

% slanted
\margin{\deleted{Si $X=\varnothing$, on prend \unsure{comme} dans} \unsure{cas}}

\deleted{\unsure{Si|donc}} $Y_0$ est un summand non vide de $Y$ tel que $f$ se factorise par $Y_1=Y\smallsetminus Y_0$, alors $f$ est \deleted{$f$ est} somme de
\[ f_{\unsure{1}}=X=X_1\longrightarrow Y_1\quad\text{et}\quad f_0:X_0=\varnothing\longrightarrow Y_0 \]
dans $\mathrm{Cat}$, donc aussi dans $\mathrm{Hot}_W$. Donc on a une somme $f=f_0\amalg f_1$ dans $\mathrm{Hot}_W$, qui est un iso., donc pour $\forall\,Z$ dans $\mathrm{Hot}_W$,

\deleted{\[ \mathrm{Hom}_W(Y_0,Z)\times\mathrm{Hom}_W(Y_1,Z) \]}
\[ \mathrm{Hom}_W(Y_0,Z)\times\mathrm{Hom}_W(Y_1,Z)
 \xrightarrow{\sim}\underbrace{\mathrm{Hom}_W(X_0,Z)}_{\substack{\parallel\\ e}}\times\mathrm{Hom}_W(X_1,Z) \]

(NB $\varnothing$ est objet initial de $\mathrm{Hot}_W$, l'étant de $\mathrm{Cat}$) i.e. le composé
\deleted{\[ Y_0(Z)\times Y_1(Z)\xrightarrow{\mathrm{pr}_2}Y_{\unsure{1}}(Z)\longrightarrow \]}
\[ Z(Y_0)\times Z(Y_1)\xrightarrow[\mathrm{pr}_2]{}Z(Y_1)\xrightarrow[Z(f_1)]{}Z(X_1) \]
est \unsure{bij.} Comme $X=X_1\ne\deleted{\unsure{\varnothing}}\varnothing$, donc aussi $Y_1\ne\varnothing$, et $Y_0\ne\varnothing$, \deleted{les \unsure{a}} les ens. $Z(Y_0),Z(Y_1),$ $Z(\unsure{\alpha_1|X_1})$ sont $\ne\varnothing$ si $Z\ne\varnothing$ ($\mathrm{Cat}$). \deleted{\unsure{Cela i}} Le composé \unsure{étant|\deleted{est} étant} bijectif, la première flèche $\mathrm{pr}_2$ est injective, ce qui implique
